\documentclass[12pt]{article}
\usepackage[utf8]{inputenc}
\usepackage[spanish]{babel}
\usepackage{amsmath}
\usepackage{amsfonts}
\usepackage{amssymb}
\usepackage{geometry}
\geometry{letterpaper, margin=2.5cm}

\begin{document}

\begin{center}
    \Large \textbf{Evaluación de Examen 1: Unidad I} \\
    \large Física para Ciencias de la Salud (005-1713) \\
    \vspace{0.5cm}
    \textbf{Estudiante:} Estefani Santoya \quad \textbf{C.I.:} 31.686.375
    \vspace{0.3cm} \\
    \large \textbf{Calificación Final: 8/10 puntos}
\end{center}

\vspace{0.5cm}

\section*{Análisis de Resultados}

\subsection*{1. Ángulo plano (Conversión a radianes)}
\textbf{Procedimiento y resultado:} Plantea una regla de tres basándose en la equivalencia $90^\circ = \frac{\pi}{2}$ rad. Desarrolla la operación multiplicando $75^\circ$ por $\pi$ y dividiendo entre $180^\circ$, obteniendo el valor decimal aproximado correcto de $0,41\pi$ rad. \\
\textbf{Veredicto:} \textbf{Correcto} (2/2 pts).

\subsection*{2. Sistemas de coordenadas (Longitud del hueso)}
\textbf{Procedimiento y resultado:} Excelente resolución. Dibuja el gráfico y ubica los puntos para formar el triángulo rectángulo. Determina con éxito las longitudes de los catetos ($a=6$, $b=8$) y aplica el Teorema de Pitágoras de manera impecable ($\sqrt{6^2 + 8^2} = \sqrt{100}$) para obtener $L = 10$. \\
\textbf{Veredicto:} \textbf{Correcto} (2/2 pts).

\subsection*{3. Vectores (Fuerza resultante)}
\textbf{Procedimiento y resultado:} Plantea correctamente la suma vectorial. Agrupa de manera adecuada las componentes $\hat{i}$ ($45 - 15 = 30$) y $\hat{j}$ ($25 + 35 = 60$), respetando la ley de los signos durante la suma aritmética, llegando al vector resultante correcto de $\vec{F} = 30\hat{i} + 60\hat{j}$ N. \\
\textbf{Veredicto:} \textbf{Correcto} (2/2 pts).

\subsection*{4. Potencias de diez (Notación científica y prefijo)}
\textbf{Procedimiento y resultado:} Expresa la medida utilizando correctamente una potencia de base diez ($12,5 \times 10^{-6}$ m). Sin embargo, la estudiante no completó la segunda parte de la pregunta, ya que omitió identificar y nombrar el prefijo del Sistema Internacional correspondiente (micrómetros, $\mu$m). \\
\textbf{Veredicto:} \textbf{Incompleto / Parcialmente correcto} (1/2 pts).

\subsection*{5. Sistemas de unidades (Conversión de densidad)}
\textbf{Procedimiento y resultado:} Intenta resolver el problema aplicando factores de conversión en cadena. Plantea correctamente la conversión de gramos a kilogramos ($1 \text{ kg} = 1000 \text{ g}$). Sin embargo, comete un \textbf{error conceptual en la conversión de volumen}, indicando que $1 \text{ m}^3 = 100 \text{ cm}^3$ en lugar del valor correcto ($1.000.000 \text{ cm}^3$). Por causa de este factor incorrecto, obtiene un resultado explícito erróneo de $0,103 \text{ kg/m}^3$. \\
\textbf{Veredicto:} \textbf{Parcialmente correcto} (Estructura de factores de conversión planteada, pero error en equivalencia de volumen) (1/2 pts).

\vspace{0.5cm}
\section*{Comentarios Generales y Sugerencias}
La estudiante demuestra un buen dominio en geometría, trigonometría básica y operaciones de vectores, manteniendo sus resoluciones claras y concisas. 
\begin{itemize}
    \item Se recomienda repasar las conversiones de unidades al cubo (volumen), recordando que al pasar de una unidad lineal a otra cúbica, el factor de conversión también debe elevarse a la potencia de tres (Ejemplo: $(100 \text{ cm})^3 = 10^6 \text{ cm}^3$).
    \item Prestar atención a todas las instrucciones de los enunciados para evitar perder puntos por respuestas incompletas (como en la Pregunta 4).
\end{itemize}

\end{document}